\section{Introduction}

Language-model agents increasingly execute quantitative research workflows that
combine data inspection, statistical estimation, strategy implementation,
portfolio construction, and executable validation. Existing quantitative
self-improvement systems primarily retain validated experiments to improve later
hypotheses, factors, or model configurations. This improves the research output,
but generally leaves the research-agent scaffold fixed.

In parallel, harness optimizers such as AHE and Meta-Harness demonstrate that
the prompts, tools, memory, middleware, and workflows surrounding a frozen model
can themselves be optimized from execution histories and task outcomes
\citep{lin2026ahe,lee2026metaharness}. We adopt this harness-level evolutionary
object but ask a domain-specific question: what feedback should guide harness
evolution for quantitative research?

Our object of adaptation is therefore the quantitative researcher, not the
strategy produced on one task: model weights remain frozen while the prompts,
tools, memory, validation, and workflow around that model may change. This
distinction also determines the evidence standard. A task-specific answer rule
is not a reusable researcher improvement even when it raises an official score.

Sparse task outcomes are often insufficient. The same wrong return, failed
property, or implausible backtest can reflect a mismatch in the task mandate,
evidence and data, quantitative representation, research operation, result
reconciliation, or artifact completion. We use these six general Research
States as an open representation rather than a fixed strategy pipeline. The
Evolver reconstructs a task-conditioned expected state and the state realized
by the Worker trajectory, maintains competing explanations of the earliest
consequential mismatch, then selects a harness component intended to mediate a
specific expected-to-observed-to-target state transition.

Many predicted transitions admit executable observations without revealing the
final answer. Changing future observations should not alter earlier decisions;
estimators should respect the declared fitting window; transaction costs should
not improve net performance; portfolio weights should satisfy the stated
exposure convention; and the submitted artifact should replay in a fresh
process. We use these task-conditioned quantitative invariants as transition
observations when applicable. Later trajectories update component experience
using actual activation, state transitions, official outcomes, and repeat or
protection evidence.

Our intended contributions are:

\begin{itemize}
    \item \Proposed{An open six-state representation that maps a quantitative
    Worker trajectory into task-conditioned expected, observed, and target
    Research States.}
    \item \Proposed{A full-harness search loop that connects a state mismatch
    to a selected component and a falsifiable state-transition observation,
    including executable quantitative invariants when applicable.}
    \item \Proposed{A candidate-level information boundary: an arm-blind
    Reviewer audits the complete Worker-visible mutation and its
    decision-changing claims before any fresh Worker evaluation.}
    \item \Proposed{A matched cumulative-campaign evaluation that separates
    information-set eligibility, component activation or prompt exposure,
    Research State correction, official improvement, stable protection-safe
    promotion, and sealed performance.}
\end{itemize}

The main comparison starts Quant-H0, Generic trajectory-guided full-harness
evolution, and QRS-guided full-harness evolution from the same frozen model and
minimal harness. Generic and QRS share the public evidence, mutation surface,
budget, Candidate Information-Set Review, and downstream evaluation gates; QRS
alone receives the State Card and state-conditioned retrieval and routing.

\Boundary{A retrospective audit found answer-rich semantic projection in all
six previously evaluated development candidates: raw diagnostics remained
Worker-hidden, but diagnostic-only predicates entered Worker-visible candidate
surfaces. Their scores remain measured development records, but they do not
establish clean harness improvement. The main comparison and sealed evaluation
therefore remain prospective.}
